% ============================================================
% A-O-08.tex -- Лекция 8 (1.8). Графы и теория чисел
% ============================================================

\documentclass[serif, aspectratio=169]{beamer}

\usepackage[T2A,T1]{fontenc}
\usepackage[english,russian]{babel}
\usepackage{XCharter}
\usepackage{amsmath}
\usepackage{amssymb}
\usepackage{graphicx}
\usepackage{tikz}
\usetikzlibrary{arrows.meta,positioning,calc,fit,overlay-beamer-styles}
\usetikzlibrary{decorations.pathmorphing}
\usetikzlibrary{braids}
\usetikzlibrary{intersections,spath3,hobby}

\usepackage{caption}
\usepackage{subcaption}
\usepackage{array}
\usepackage{booktabs}
\usepackage{mathtools}
\usepackage{dynkin-diagrams}
%\usepackage{fontawesome5}

\newcommand{\sagemark}{%
  \textcolor{green!40!black}{\textbf{\texttt{sage:...}}}%
}

\usetheme{Math}
\usepackage{mymacros}
% \colorscheme{green}
% \colorscheme{terracotta}

% ------------------------------------------------------------
% Title data
% ------------------------------------------------------------
\author{Павел Владимирович Снурницын}
\title{Лекция 8 (1.8)\\ Графы и теория чисел}
\titlecontext{Инструменты машинно-вспомогательной арифметики\\ с/к, осень 2026}
\institute{к.ф.-м.н., с.н.с., ИБ ВМК}
\date{\small 17 ноября 2026 г.}

\setfooterauthor{П.В. Снурницын}
\setfootertopic{1.8 Графы и теория чисел}
\setfootercontext{Спецкурс A-O 2026}

\newif\ifanimated
\newcommand{\animateon}{\beamerdefaultoverlayspecification{<+-| alert@+>}\animatedtrue}
\newcommand{\animateoff}{\beamerdefaultoverlayspecification{}\animatedfalse}
\newcommand{\steponhelper}[2]{% #1 = свойство цвета (text/draw), #2 = номер шага
  \ifanimated
    \edef\stepon@tmp{alt=<#2>{#1=ThemeAlert}{alt=<#2->{#1=black}{invisible}}}%
  \else
    \edef\stepon@tmp{#1=black}%
  \fi
  \expandafter\pgfkeysalso\expandafter{\stepon@tmp}%
}
\tikzset{
  step on/.code={\steponhelper{text}{#1}},
  step draw on/.code={\steponhelper{draw}{#1}},
  step fill on/.code={\steponhelper{fill}{#1}},
}
% \animateon
\animateoff

\begin{document}

% ================================================================
\begin{frame}
    \titlepage
\end{frame}
\note{}

% ================================================================
\begin{frame}{В прошлый раз}
    \begin{itemize}
        \item $p$-адические числа: $\mathbb{Z}_p$, $\mathbb{Q}_p$;
        \item Также вспоминали:
        \begin{itemize}[<.->]
            \item конечные поля (и их расширения),
            \item кватернионные алгебры,
            \item модулярные формы,
            \item эллиптические кривые,
            \item дзета-функции.
        \end{itemize}
    \end{itemize}
    \note{
        \begin{itemize}
            \item ---
        \end{itemize}
    }
\end{frame}

% ============================================================
\section{Арифметические графы}
% ============================================================

% ================================================================
\begin{frame}{Арифметические графы}
    \begin{itemize}
        \item Граф $G=(V,E)$: множество вершин $V=\{v\}$ и рёбер $E=\{\{v_1,v_2\}\}$;
        \item Рёбра могут быть ориентированными, кратными или петлями;
        \item Идея:
        \begin{itemize}
            \item $V$ --- арифметические (математические) объекты,
            \item $E$ --- отношения или преобразования между ними,
        \end{itemize}
        \item Какие примеры таких графов уже видели ранее:
        \begin{center}
            \begin{actionenv}
            \begin{minipage}[t]{0.47\linewidth}
                \centering
                \begin{tikzpicture}[
                    scale=0.84,
                    every node/.style={font=\footnotesize,inner sep=2pt}
                ]
                    \path[use as bounding box]
                        (-2.65,-1.45) rectangle (2.65,1.60);

                    % A_3
                    \begin{scope}[shift={(-1.25,0.85)}]
                        \draw[thick] (-0.60,0)--(0,0)--(0.60,0);
                        \foreach \x in {-0.60,0,0.60}
                            \fill (\x,0) circle[radius=2pt];
                        \node at (0,-0.30) {$A_3$};
                    \end{scope}

                    % D_4
                    \begin{scope}[shift={(1.25,0.85)}]
                        \draw[thick]
                            (-0.60,0)--(0,0)--(0.60,0)
                            (0,0)--(0,0.60);
                        \foreach \x/\y in {-0.60/0,0/0,0.60/0,0/0.60}
                            \fill (\x,\y) circle[radius=2pt];
                        \node at (0,-0.30) {$D_4$};
                    \end{scope}

                    % E_8
                    \begin{scope}[yshift=-0.75cm]
                        \draw[thick]
                            (-1.80,0)--(1.80,0)
                            (-0.60,0)--(-0.60,0.60);
                        \foreach \x in {-1.80,-1.20,-0.60,0,0.60,1.20,1.80}
                            \fill (\x,0) circle[radius=2pt];
                        \fill (-0.60,0.60) circle[radius=2pt];
                        \node at (0,-0.30) {$E_8$};
                    \end{scope}
                \end{tikzpicture}\\[-0.5ex]
                {\scriptsize Диаграмма Дынкина $A_3\oplus D_4\oplus E_8$}
            \end{minipage}
            \end{actionenv}
            \begin{actionenv}
            \begin{minipage}[t]{0.47\linewidth}
                \centering
                \begin{tikzpicture}[
                    scale=0.84,
                    every node/.style={font=\footnotesize,inner sep=2pt, outer sep=2pt}
                ]
                    \path[use as bounding box]
                        (-2.65,-1.45) rectangle (2.65,1.60);

                    % Вершины двоичного дерева
                    \coordinate (r) at (0,1.12);
                    \coordinate (a0) at (-1.20,0.42);
                    \coordinate (a1) at ( 1.20,0.42);

                    \foreach \k/\x in {0/-1.80,1/-0.60,2/0.60,3/1.80}
                        \coordinate (b\k) at (\x,-0.28);

                    \foreach \k in {0,...,7}
                        \coordinate (c\k) at ({-2.10+0.60*\k},-0.98);

                    % Рёбра
                    \draw[thick]
                        (r)--(a0) (r)--(a1)
                        (a0)--(b0) (a0)--(b1)
                        (a1)--(b2) (a1)--(b3);

                    \foreach \u/\v in {
                        0/0,0/1,1/2,1/3,2/4,2/5,3/6,3/7}
                        \draw[thick] (b\u)--(c\v);

                    % Вершины
                    \foreach \v in {
                        r,a0,a1,b0,b1,b2,b3,
                        c0,c1,c2,c3,c4,c5,c6,c7}
                        \fill (\v) circle[radius=2pt];

                    % Представители классов вычетов
                    \node[above] at (r) {$\mathbb Z_2$};

                    \node[above left]  at (a0) {$0$};
                    \node[above right] at (a1) {$1$};

                    \node[above left]  at (b0) {$0$};
                    \node[above right] at (b1) {$2$};
                    \node[above left]  at (b2) {$1$};
                    \node[above right] at (b3) {$3$};

                    \foreach \k/\val in {0/0,1/4,2/2,3/6,4/1,5/5,6/3,7/7}
                        \node[below] at (c\k) {$\val$};
                \end{tikzpicture}\\[-0.5ex]
                {\scriptsize Приближения в $\mathbb Z_2$}
            \end{minipage}
            \end{actionenv}
        \end{center}
        \vspace{1ex}
        \item Можно изучать св-ва математических объектов методами теории графов;
        \item А можно строить графы с «хорошими» свойствами по этим объектам.
    \end{itemize}
    \note{
        \begin{itemize}
            \item ---
        \end{itemize}
    }
\end{frame}

% ================================================================
\begin{frame}{Башни расширений полей}
    \begin{itemize}
        \item Пусть $p$ --- простое, $\mathbb{F}_p=\mathbb{Z}/p\mathbb{Z}$ --- конечное поле;
        \item Теорема: $\mathbb{F}_{p^m}\subseteq\mathbb{F}_{p^n}$ $\Leftrightarrow$ $m\mid n$; $\mathbb{F}_{p^n}$ --- расширение $\mathbb{F}_{p^m}$;
        \item Пример: $\mathbb{F}_{p}\subset\mathbb{F}_{p^2}\subset\mathbb{F}_{p^4}\subset\mathbb{F}_{p^8}\subset\dots$;
        \item Граф:
        \begin{itemize}
            \item вершины --- поля, 
            \item рёбра --- (непосредственные) включения;
        \end{itemize}
        \item Примеры:
        \vspace{-3ex}
        \begin{center}
            \begin{actionenv}
            \begin{minipage}[t]{0.4\linewidth}
                \centering
                \begin{tikzpicture}[
                    node distance=6mm,
                    fld/.style={inner sep=2pt, font=\scriptsize},
                    >={Stealth}
                ]
                    \node[fld] (F1) {$\mathbb{F}_{p}$};
                    \node[fld, above=of F1] (F2) {$\mathbb{F}_{p^2}$};
                    \node[fld, above=of F2] (F4) {$\mathbb{F}_{p^4}$};

                    \draw[thick] (F1) -- (F2);
                    \draw[thick] (F2) -- (F4);
                \end{tikzpicture}\\
                {\scriptsize Башня расширений}
            \end{minipage}
            \end{actionenv}
            \begin{actionenv}
            \begin{minipage}[t]{0.4\linewidth}
                \centering
                \begin{tikzpicture}[
                    fld/.style={inner sep=2pt, font=\scriptsize}
                ]
                    \node[fld] (F1)  at (0,0)    {$\mathbb{F}_{p}$};
                    \node[fld] (F2)  at (-1.35,1.0) {$\mathbb{F}_{p^2}$};
                    \node[fld] (F3)  at ( 1.35,1.0) {$\mathbb{F}_{p^3}$};
                    \node[fld] (F4)  at (-1.35,2.0) {$\mathbb{F}_{p^4}$};
                    \node[fld] (F6)  at ( 1.35,2.0) {$\mathbb{F}_{p^6}$};
                    \node[fld] (F12) at (0,3.0)    {$\mathbb{F}_{p^{12}}$};

                    \draw[thick] (F1) -- (F2);
                    \draw[thick] (F1) -- (F3);
                    \draw[thick] (F2) -- (F4);
                    \draw[thick] (F2) -- (F6);
                    \draw[thick] (F3) -- (F6);
                    \draw[thick] (F4) -- (F12);
                    \draw[thick] (F6) -- (F12);
                \end{tikzpicture}\\
                {\scriptsize Подполя $\mathbb{F}_{p^{12}}$}
            \end{minipage}
            \end{actionenv}
        \end{center}
    \end{itemize}
    \note{
        \begin{itemize}
            \item ---
        \end{itemize}
    }
\end{frame}

% ================================================================
\begin{frame}{Граф Пэли}
    \begin{itemize}
        \item Пусть $p\equiv1\pmod4$ --- простое;
        \item Граф Пэли $P(p)$: вершины $V=\mathbb{F}_p$, рёбра $(a,b)\in E$ $\Leftrightarrow$ $a\neq b$, $a-b$ квадрат в $\mathbb{F}_p$;
        \item $-1$ --- квадрат при $p\equiv1\pmod4$, поэтому отношение симметрично;
        \item Кстати, $a\in\mathbb{Z}$ такое, что разрешимо сравнение $a\equiv x^2 \pmod{n}$, называется квадратичным вычетом $\bmod\, n$ (если неразрешимо, то невычетом);
        \item Примеры:
        \vspace{-2ex}
        \begin{center}
            \begin{actionenv}
            \begin{minipage}[t]{0.45\linewidth}
                \centering
                \begin{tikzpicture}[
                    scale=0.8,
                    every node/.style={font=\footnotesize, inner sep=1pt}
                ]
                    \path[use as bounding box] (-2,-2) rectangle (2,2);

                    \foreach \k in {0,...,4}{
                        \coordinate (v\k) at ({90-360*\k/5}:1.35);
                        \node at ({90-360*\k/5}:1.72) {$\k$};
                    }

                    \draw[thick] (v0) -- (v1) -- (v2) 
                        -- (v3) -- (v4) -- (v0);

                    \foreach \k in {0,...,4}
                        \fill (v\k) circle[radius=2pt];
                \end{tikzpicture}\\[-0.5ex]
                {\scriptsize $P(5)$: ненулевые квадраты $\{1,4\}$}
            \end{minipage}
            \end{actionenv}
            \begin{actionenv}
            \begin{minipage}[t]{0.45\linewidth}
                \centering
                \begin{tikzpicture}[
                    scale=0.8,
                    every node/.style={font=\footnotesize, inner sep=1pt}
                ]
                    \path[use as bounding box] (-2,-2) rectangle (2,2);

                    \foreach \k in {0,...,12}{
                        \coordinate (v\k) at ({90-360*\k/13}:1.35);
                        \node at ({90-360*\k/13}:1.72) {$\k$};
                    }

                    \foreach \k in {0,...,12}{
                        \foreach \d in {1,3,4}{
                            \pgfmathtruncatemacro{\j}{mod(\k+\d,13)}
                            \draw[thick] (v\k) -- (v\j);
                        }
                    }

                    \foreach \k in {0,...,12}
                        \fill (v\k) circle[radius=2pt];
                \end{tikzpicture}\\[-0.5ex]
                {\scriptsize $P(13)$: ненулевые квадраты $\{1,3,4,9,10,12\}$}
            \end{minipage}
            \end{actionenv}
        \end{center}
    \end{itemize}
    \note{
        \begin{itemize}
            \item ---
        \end{itemize}
    }
\end{frame}

% ================================================================
\begin{frame}{Графы Кэли}
    \begin{itemize}
        \item Пусть $H$ --- конечная группа, $S\subset H$ --- множество образующих;
        \item Граф Кэли $\operatorname{Cay}(H,S)$: вершины $V=H$, рёбра $(h_1, h_2)\in E$ $\Leftrightarrow$ $\exists\, s\in S$: $h_2=h_1 s$;
        \item При $S=S^{-1}$ граф можно считать неориентированным; при $e\notin S$ нет петель;

        \item
        \begin{minipage}[t]{0.68\linewidth}
            \raggedright
            \strut Пример: $H=\mathbb{Z}/n\mathbb{Z}$, $S=\{\pm1\}$ при $n\ge3$ ---
            цикл; все вершины имеют одну и ту же локальную структуру;
        \end{minipage}\hfill
        \begin{minipage}[t]{0.30\linewidth}
            \centering
            \begin{tikzpicture}[
                baseline={([yshift={-\ht\strutbox}]current bounding box.north)},
                scale=0.7,
                every node/.style={font=\footnotesize, inner sep=1pt}
            ]
                \foreach \k in {0,...,7}{
                    \coordinate (v\k) at ({90-45*\k}:1.2);
                    \node at ({90-45*\k}:1.48) {$\k$};
                }

                \draw[thick] (v0) -- (v1) -- (v2) -- (v3)
                    -- (v4) -- (v5) -- (v6) -- (v7) -- (v0);

                \foreach \k in {0,...,7}
                    \fill (v\k) circle[radius=2pt];
            \end{tikzpicture}
        \end{minipage}

        \item
        \begin{minipage}[t]{0.68\linewidth}
            \raggedright
            \strut Пример: группа симметрий квадрата
            $D_4=\langle r,s\mid r^4=s^2=1,\ srs=r^{-1}\rangle$,
            $S=\{r,r^{-1},s\}$.
        \end{minipage}\hfill
        \begin{minipage}[t]{0.30\linewidth}
            \centering
            \begin{tikzpicture}[
                baseline={([yshift={-\ht\strutbox}]current bounding box.north)},
                scale=0.7,
                every node/.style={font=\footnotesize, inner sep=1pt,  outer sep=2pt}
            ]
                % Внутренний квадрат
                \coordinate (e)   at (-0.60, 0.60);
                \coordinate (r)   at ( 0.60, 0.60);
                \coordinate (r2)  at ( 0.60,-0.60);
                \coordinate (r3)  at (-0.60,-0.60);

                % Внешний квадрат
                \coordinate (s)   at (-1.35, 1.35);
                \coordinate (rs)  at ( 1.35, 1.35);
                \coordinate (r2s) at ( 1.35,-1.35);
                \coordinate (r3s) at (-1.35,-1.35);

                % Рёбра
                \draw[thick]
                    (e) -- (r) -- (r2) -- (r3) -- (e)
                    (s) -- (rs) -- (r2s) -- (r3s) -- (s)
                    (e) -- (s)
                    (r) -- (rs)
                    (r2) -- (r2s)
                    (r3) -- (r3s);

                % Вершины
                \foreach \v in {e,r,r2,r3,s,rs,r2s,r3s}
                    \fill (\v) circle[radius=2pt];

                % Обозначения вершин
                \node[below right] at (e) {$1$};
                \node[below left]  at (r) {$r$};
                \node[above left]  at (r2) {$r^2$};
                \node[above right] at (r3) {$r^3$};

                \node[above left]  at (s) {$s$};
                \node[above right] at (rs) {$rs$};
                \node[below right] at (r2s) {$r^2s$};
                \node[below left]  at (r3s) {$r^3s$};
            \end{tikzpicture}
        \end{minipage}
    \end{itemize}
    \note{
        \begin{itemize}
            \item ---
        \end{itemize}
    }
\end{frame}

% ================================================================
\begin{frame}{Некоторые понятия из теории графов}
    \begin{itemize}
        \item Простой граф $\leftrightharpoons$ неориентированный граф без петель и кратных рёбер;
        \item Степень вершины $\deg(v)$ $\leftrightharpoons$ число инцидентных ей рёбер (петли --- дважды);
        \item Лемма (о сумме степеней): $\sum_{v\in V}\deg(v)=2|E|$;
        \item $d$-регулярный граф $\leftrightharpoons$ $\deg(v)=d$ для всех $v\in V$; $d|V|=2|E|$;
        \item Пример: граф Пэли --- $(p-1)/2$-регулярный ($(p-1)/2$ квадратов $\neq 0$);
        \item Маршрут $\leftrightharpoons$ посл-сть смежных вершин, длина --- число пройденных рёбер;
        \item Путь $\leftrightharpoons$ маршрут без повторения вершин;
        \item Цикл $\leftrightharpoons$ замкнутый маршрут без повторения (кроме начала и конца);
        \item Связный граф $\leftrightharpoons$ между любыми двумя вершинами существует путь;
        \item Расстояние $d_G(u,v)$ $\leftrightharpoons$ длина кратчайшего пути между $u$ и $v$;
        \item Диаметр связного графа $\leftrightharpoons$ $\operatorname{diam}(G)=\sup_{u,v\in V}d_G(u,v)$;
        \item Дерево $\leftrightharpoons$ связный граф без циклов;
        \item Двудольный граф $\leftrightharpoons$ $V=V_1\cup V_2$, $V_1\cap V_2=\varnothing$ и $\forall \{v_1,v_2\}\in E$: $v_1\in V_1$, $v_2\in V_2$.
    \end{itemize}
    \note{
        \begin{itemize}
            \item ---
        \end{itemize}
    }
\end{frame}

% ================================================================
\begin{frame}{Матрица смежности}
    \begin{itemize}
        \item Для графа без кратных рёбер $n$ вершинами матрица смежности:
        \[A=(a_{ij})_{i,j=1}^n,\qquad
        a_{ij}=\begin{cases}
            1,&(v_i,v_j)\in E,\\
            0,&(v_i,v_j)\notin E;
        \end{cases}\]
        \item $A=A^{\mathsf{T}}$ для неориентированного $G$; $a_{ii}=1$, если есть петля в $v_i$;
        \item $(A^k)_{ij}$ --- число маршрутов длины $k$ из $v_i$ в $v_j$ (с повторениями);
        \item $\operatorname{tr}(A^k)$ --- число замкнутых маршрутов с отмеченным началом;
        \item Пример: $C_4=\operatorname{Cay}(\mathbb{Z}/4\mathbb{Z},\{\pm1\})$, $v_i=i-1$,
            \begin{center}
            \begin{tikzpicture}[
                baseline=(current bounding box.center),
                >={Stealth[length=2mm]},
                every node/.style={font=\footnotesize, inner sep=1pt, outer sep=2pt}
            ]
                \ifanimated
                    \tikzset{
                        step1/.style={visible on=<.->, alt=<.>{color=ThemeAlert}{}},
                        step2/.style={visible on=<+->, alt=<.>{color=ThemeAlert}{}},
                        step3/.style={visible on=<+->, alt=<.>{color=ThemeAlert}{}},
                        step4/.style={visible on=<+->, alt=<.>{color=ThemeAlert}{}},
                        step5/.style={visible on=<+->, alt=<.>{color=ThemeAlert}{}}
                    }
                \else
                    \tikzset{
                        step1/.style={},
                        step2/.style={},
                        step3/.style={},
                        step4/.style={},
                        step5/.style={}
                    }
                \fi

                \path[use as bounding box] (-1.3,-2.2) rectangle (9.6,1.3);

                % --- Граф C_4 ---
                \coordinate (v1) at (-0.56, 0.56);
                \coordinate (v2) at ( 0.56, 0.56);
                \coordinate (v3) at ( 0.56,-0.56);
                \coordinate (v4) at (-0.56,-0.56);

                \begin{scope}[step1]
                    \draw[thick] (v1) -- (v2) -- (v3) -- (v4) -- (v1);

                    \foreach \i in {1,...,4}
                        \fill (v\i) circle[radius=2pt];

                    \node[above left]  at (v1) {$v_1$};
                    \node[above right] at (v2) {$v_2$};
                    \node[below right] at (v3) {$v_3$};
                    \node[below left]  at (v4) {$v_4$};
                \end{scope}

                % --- Матрица A ---
                \node[anchor=west, step2] at (1.55,0) {$\displaystyle
                    A=\begin{pmatrix}
                        0&1&0&1\\
                        1&0&1&0\\
                        0&1&0&1\\
                        1&0&1&0
                    \end{pmatrix}
                $};

                % --- Матрица A^2 ---
                \node[anchor=west, step3] at (4.8,0) {$\displaystyle
                    A^2=\begin{pmatrix}
                        2&0&2&0\\
                        0&2&0&2\\
                        2&0&2&0\\
                        0&2&0&2
                    \end{pmatrix}
                $};

                % --- (A^2)_{13} и соответствующие пути ---
                \begin{scope}[step4]
                    % Путь v1 -> v2 -> v3
                    \draw[->,thick,draw=green!60!black]
                        ($(v1)!0.18!(v2)+(0,-0.14)$) --
                        ($(v1)!0.82!(v2)+(0,-0.14)$);
                    \draw[->,thick,draw=green!60!black]
                        ($(v2)!0.18!(v3)+(-0.14,0)$) --
                        ($(v2)!0.82!(v3)+(-0.14,0)$);

                    % Путь v1 -> v4 -> v3
                    \draw[->,thick,draw=red]
                        ($(v1)!0.18!(v4)+(0.14,0)$) --
                        ($(v1)!0.82!(v4)+(0.14,0)$);
                    \draw[->,thick,draw=red]
                        ($(v4)!0.18!(v3)+(0,0.14)$) --
                        ($(v4)!0.82!(v3)+(0,0.14)$);

                    \node[anchor=north west, align=left] at (-1.0,-1.20) {$\displaystyle
                        (A^2)_{13}=2:\quad
                        \textcolor{green!60!black}{v_1\to v_2\to v_3},\quad
                        \textcolor{red}{v_1\to v_4\to v_3};
                    $};
                \end{scope}

                % --- След ---
                \node[anchor=north west, step5] at (6.3,-1.20) {$\displaystyle
                    \operatorname{tr}(A^2)=8.
                $};
            \end{tikzpicture}
            \end{center}
    \end{itemize}
    \note{
        \begin{itemize}
            \item При кратных рёбрах соответствующие элементы учитывают их число $a_{ij}$ --- число рёбер;
        \end{itemize}
    }
\end{frame}

% ================================================================
\begin{frame}{Спектр регулярного графа}
    \begin{itemize}
        \item Пусть $G$ --- конечный простой граф, $A=A^{\mathsf T}$ --- его матрица смежности;
        \item Число $\lambda$ --- собственное значение $A$, если $A\mathbf{x}=\lambda \mathbf{x}$ при некотором $\mathbf{x}\ne\mathbf{0}$;\\
        для симметричной $A$ все $\lambda$ действительны;
        \item Спектр графа $\leftrightharpoons$ собственные значения $A$: $\lambda_1\ge\lambda_2\ge\cdots\ge\lambda_n$;
        \item Лемма: Для неориентированного $d$-регулярного графа без петель кратность собственного значения $d$ равна числу компонент связности.
        \item Лемма: Для связного $d$-регулярного графа $\lambda_1=d$,\quad $-d\le\lambda_i<d$, $i\ge2$;\\
        $-d$ --- собственное значение ровно тогда, когда граф двудольный;
        \item $\lambda=\pm d$ будем называть тривиальными собственными значениями;
        \item Спектральный зазор $\leftrightharpoons$ $\lambda_1-\lambda_2$; для связного $d$-регулярного: $d-\lambda_2$;
        \item Т.о. спектр графа --- инвариант, отражающий связность и некоторые другие характеристики;
        \item Пример: спектр цикла $C_4$ --- $(2,0,0,-2)$, зазор --- $2$.
    \end{itemize}
    \note{
        \begin{itemize}
            \item ---
        \end{itemize}
    }
\end{frame}

% ============================================================
\section{Графы-экспандеры}
% ============================================================

% ================================================================
\begin{frame}{«Плохие» и «хорошие» графы}
    \begin{center}
        % Первая строка: плохие примеры
        \begin{actionenv}
        \begin{minipage}[t]{0.31\linewidth}
            \centering
            
            \begin{tikzpicture}[scale=0.78]
                \path[use as bounding box] (-1.85,-1.45) rectangle (1.85,1.45);

                \foreach \k in {0,...,9}
                    \coordinate (v\k) at ({90-36*\k}:1.35);

                \draw[thick] (v0)--(v1)--(v2)--(v3)--(v4)
                    --(v5)--(v6)--(v7)--(v8)--(v9)--(v0);

                \foreach \k in {0,...,9}
                    \fill (v\k) circle[radius=2pt];
            \end{tikzpicture}\par
            {\scriptsize Цикл $C_{10}$}\par
            {\scriptsize Мало рёбер, большие расстояния}
        \end{minipage}
        \end{actionenv}
        \begin{actionenv}
        \begin{minipage}[t]{0.31\linewidth}
            \centering
            \begin{tikzpicture}[scale=0.78]
                \path[use as bounding box] (-1.85,-1.45) rectangle (1.85,1.45);

                \begin{scope}[xshift=-1.05cm]
                    \foreach \k in {0,...,4}
                        \coordinate (l\k) at ({72*\k}:0.70);
                \end{scope}
                \begin{scope}[xshift=1.05cm]
                    \foreach \k in {0,...,4}
                        \coordinate (r\k) at ({180+72*\k}:0.70);
                \end{scope}

                \foreach \i in {0,...,4}{
                    \foreach \j in {0,...,4}{
                        \ifnum\i<\j
                            \draw[thick] (l\i)--(l\j);
                            \draw[thick] (r\i)--(r\j);
                        \fi
                    }
                }

                \draw[thick] (l0)--(r0);

                \foreach \k in {0,...,4}{
                    \fill (l\k) circle[radius=2pt];
                    \fill (r\k) circle[radius=2pt];
                }
            \end{tikzpicture}\par
            {\scriptsize Два кластера и мост}\par
            {\scriptsize Узкое место}
        \end{minipage}
        \end{actionenv}
        \begin{actionenv}
        \begin{minipage}[t]{0.31\linewidth}
            \centering
            \begin{tikzpicture}[scale=0.78]
                \path[use as bounding box] (-1.85,-1.45) rectangle (1.85,1.45);

                \foreach \k in {0,...,7}
                    \coordinate (v\k) at ({90-45*\k}:1.35);

                \foreach \i in {0,...,7}{
                    \foreach \j in {0,...,7}{
                        \ifnum\i<\j
                            \draw[semithick] (v\i)--(v\j);
                        \fi
                    }
                }

                \foreach \k in {0,...,7}
                    \fill (v\k) circle[radius=2pt];
            \end{tikzpicture}\par
            {\scriptsize Полный граф $K_8$}\par
            {\scriptsize Слишком много рёбер}
        \end{minipage}
        \end{actionenv}
        \vspace{2ex}

        % Вторая строка: хорошие примеры
        \begin{minipage}[t]{0.32\linewidth}
            \vspace{-13ex}
            \begin{itemize}
                \item Т.о. «хороший» граф:
                \begin{itemize}
                    \item много узлов,
                    \item мало рёбер,
                    \item хорошо свзяный,
                    \item без узких мест;
                \end{itemize}
            \end{itemize}
        \end{minipage}
        \begin{actionenv}
        \begin{minipage}[t]{0.32\linewidth}
            \centering
            \begin{tikzpicture}[scale=0.82]
                \path[use as bounding box] (-1.85,-1.45) rectangle (1.85,1.45);

                \foreach \k in {0,...,4}{
                    \coordinate (v\k) at ({90-72*\k}:1.35);
                    \coordinate (u\k) at ({90-72*\k}:0.64);
                }

                \draw[thick]
                    (v0)--(v1)--(v2)--(v3)--(v4)--(v0)
                    (u0)--(u2)--(u4)--(u1)--(u3)--(u0);

                \foreach \k in {0,...,4}
                    \draw[thick] (v\k)--(u\k);

                \foreach \k in {0,...,4}{
                    \fill (v\k) circle[radius=2pt];
                    \fill (u\k) circle[radius=2pt];
                }
            \end{tikzpicture}\par
            {\scriptsize Граф Петерсена}
        \end{minipage}
        \end{actionenv}
        \begin{actionenv}
        \begin{minipage}[t]{0.32\linewidth}
            \centering
            \begin{tikzpicture}[scale=0.82]
                \path[use as bounding box] (-2.0,-1.55) rectangle (2.0,1.55);

                \foreach \k in {0,...,13}
                    \coordinate (v\k) at ({90-360*\k/14}:1.35);

                % цикл C_14
                \draw[thick]
                    (v0)--(v1)--(v2)--(v3)--(v4)--(v5)--(v6)--(v7)
                    --(v8)--(v9)--(v10)--(v11)--(v12)--(v13)--(v0);

                % хорды графа Хивуда: {2k,2k+5} mod 14
                \draw[thick] (v0)--(v5);
                \draw[thick] (v2)--(v7);
                \draw[thick] (v4)--(v9);
                \draw[thick] (v6)--(v11);
                \draw[thick] (v8)--(v13);
                \draw[thick] (v10)--(v1);
                \draw[thick] (v12)--(v3);

                \foreach \k in {0,...,13}
                    \fill (v\k) circle[radius=2pt];
            \end{tikzpicture}\par
            {\scriptsize Граф Хивуда}
        \end{minipage}
        \end{actionenv}
    \end{center}

    \note{
        \begin{itemize}
            \item Граф Хивуда можно получить как граф инцидентности плоскости Фано: вершины --- 7 точек и 7 прямых, ребро означает принадлежность точки прямой. Эквивалентно: точки --- одномерные подпространства $\mathbb{F}_2^3$, прямые --- двумерные подпространства, рёбра --- включения.
        \end{itemize}
    }
\end{frame}

% ================================================================
\begin{frame}{Константа Чигера}
    \begin{itemize}
        \item Для $U\subset V(G)$ обозначим через $\partial U$ множество рёбер между $U$ и $V\setminus U$;
        \item Константа Чигера $\leftrightharpoons$ $h(G)=\min\limits_{0<|U|\le |V|/2}|\partial U|/|U|$;
        \item Примеры:
        \vspace{-2ex}
        \begin{center}
            \begin{actionenv}
            \begin{minipage}[t]{0.31\linewidth}
                \centering
                \begin{tikzpicture}[scale=0.82]
                    \path[use as bounding box] (-1.95,-1.75) rectangle (1.95,1.75);

                    % Вершины цикла
                    \foreach \k in {0,...,9}
                        \coordinate (v\k) at ({90-36*\k}:1.18);

                    % Обычные рёбра
                    \draw[thick]
                        (v0)--(v1)--(v2)--(v3)--(v4)
                        (v5)--(v6)--(v7)--(v8)--(v9);

                    % Граница \partial U
                    \draw[red,thick] (v9)--(v0) (v4)--(v5);

                    % Выделенное множество U
                    \draw[red,dashed]
                        (-72:1.48)
                        arc[start angle=-72,end angle=108,radius=1.48]
                        -- cycle;

                    \node[red] at (0.7,0.3) {\scriptsize $U$};

                    % Вершины
                    \foreach \k in {0,...,9}
                        \fill (v\k) circle[radius=2pt];
                \end{tikzpicture}\\[-2ex]
                {\scriptsize $|U|=5,\ |\partial U|=2,\ h(G)=2/5$}
            \end{minipage}
            \end{actionenv}
            \begin{actionenv}
            \begin{minipage}[t]{0.31\linewidth}
                \centering
                \begin{tikzpicture}[scale=0.82]
                    \path[use as bounding box] (-2.05,-1.75) rectangle (2.05,1.75);

                    % Левый K5
                    \begin{scope}[xshift=-1.10cm]
                        \foreach \k in {0,...,4}
                            \coordinate (l\k) at ({90-72*\k}:0.72);
                    \end{scope}

                    % Правый K5
                    \begin{scope}[xshift=1.10cm]
                        \foreach \k in {0,...,4}
                            \coordinate (r\k) at ({90-72*\k}:0.72);
                    \end{scope}

                    % Рёбра K5
                    \foreach \i in {0,...,4}{
                        \foreach \j in {0,...,4}{
                            \ifnum\i<\j
                                \draw[thick] (l\i)--(l\j);
                                \draw[thick] (r\i)--(r\j);
                            \fi
                        }
                    }

                    % Мост = граница
                    \draw[red,thick] (l1)--(r4);

                    % Выделенное множество U
                    \draw[red,dashed,rounded corners=8pt]
                        (-2.02,1.05) rectangle (-0.18,-1.05);
                    \node[red] at (-0.5,0.7) {\scriptsize $U$};

                    % Вершины
                    \foreach \k in {0,...,4}{
                        \fill (l\k) circle[radius=2pt];
                        \fill (r\k) circle[radius=2pt];
                    }
                \end{tikzpicture}\\[-2ex]
                {\scriptsize $|U|=5,\ |\partial U|=1,\ h(G)=1/5$}
            \end{minipage}
            \end{actionenv}
            \begin{actionenv}
            \begin{minipage}[t]{0.31\linewidth}
                \centering
                \begin{tikzpicture}[scale=0.82]
                    \path[use as bounding box] (-1.95,-1.75) rectangle (1.95,1.75);

                    % Внешние и внутренние вершины
                    \foreach \k in {0,...,4}{
                        \coordinate (v\k) at ({90-72*\k}:1.20);
                        \coordinate (u\k) at ({90-72*\k}:0.58);
                    }

                    % Обычные рёбра графа Петерсена
                    \draw[thick]
                        (v0)--(v1)--(v2)--(v3)--(v4)--(v0)
                        (u0)--(u2)--(u4)--(u1)--(u3)--(u0);

                    % Граница \partial U
                    \foreach \k in {0,...,4}
                        \draw[red,thick] (v\k)--(u\k);

                    % Выделенное множество U
                    \draw[red,dashed] (0,0) circle[radius=0.8];
                    \node[red] at (0.35,0.45) {\scriptsize $U$};

                    % Вершины
                    \foreach \k in {0,...,4}{
                        \fill (v\k) circle[radius=2pt];
                        \fill (u\k) circle[radius=2pt];
                    }
                \end{tikzpicture}\\[-2ex]
                {\scriptsize $|U|=5,\ |\partial U|=5,\ h(G)=1$}
            \end{minipage}
            \end{actionenv}
        \end{center}
        \vspace{1ex}
        \item Интерпретация: чем больше $h(G)$, тем труднее отделить существенную часть графа малым числом рёбер;
        \item Это и будет характеристикой хороших графов: чем больше константа Чигера, тем лучше «расширение» графа.
    \end{itemize}
    \note{
        \begin{itemize}
            \item ---
        \end{itemize}
    }
\end{frame}

% ================================================================
\begin{frame}{Спектральный зазор и расширение}
    \begin{itemize}
        \item Пусть $G$ --- связный $d$-регулярный граф ($d-\lambda_2$ --- спектральный зазор);
        \item Теорема (неравенства Чигера):
        \[\frac{d-\lambda_2}{2}\le h(G)\le\sqrt{2d(d-\lambda_2)}.\]
        \item Интерпретация: большой зазор $\Rightarrow$ хорошее расширение;
        \item Пример: Граф Пэли 
        \begin{itemize}
            \item вершины $a\in\mathbb{F}_p$; $a\sim b$ при $a\ne b$ и $a-b=x^2$ для некоторого $x\in\mathbb{F}_p^*$;
            \item арифметическое описание позволяет найти спектр: $\frac{p-1}{2}$, $\frac{-1+\sqrt p}{2}$, $\frac{-1-\sqrt p}{2}$;
            % \item Последние два значения имеют кратность $(p-1)/2$ каждое и связаны с квадратичными суммами Гаусса;
            % \item Суммы характеров $\longrightarrow$ оценки спектра $\longrightarrow$ псевдослучайные свойства;
            \item и по спектру оценить $h(P(p))$:
            \[
                \frac{p-\sqrt p}{4}\le h(P(p))\le
                \sqrt{\frac{(p-1)(p-\sqrt p)}{2}}.
            \]
            \item Нижняя граница порядка $p$: даже большие части графа связаны многими рёбрами с остальными.
        \end{itemize}
        
    \end{itemize}

    \note{
        \begin{itemize}
            \item ---
        \end{itemize}
    }
\end{frame}

% ================================================================
\begin{frame}{Семейства расширяющихся графов (экспандеры)}
    \begin{itemize}
        \item У графов Пэли степень вершин $d=(p-1)/2$ растёт с $p$;
        \item Хотим семейство графов $G_n$ с $|V(G_n)|\to\infty$, но фиксированной степенью $d$;
        \item Но «хороших» (например, для семейства циклов $C_n$: $h(C_n)\to0$, $n\to\infty$);
        \item Семейство экспандеров $\leftrightharpoons$ существует $\varepsilon>0$ такое, что $h(G_n)\ge\varepsilon$ для всех $n$;
        \item Теорема (Алон--Боппана): для семейства связных $d$-регулярных графов с растущим числом вершин ($d\ge3$):
        \[\liminf_{n\to\infty}\lambda_2(G_n)\ge2\sqrt{d-1}.\]
        \item Граф называется графом Рамануджана, если каждое нетривиальное собственное значение  удовлетворяет
        \[|\lambda|\le2\sqrt{d-1},\quad \lambda\neq \pm d;\]
        % \item Это оптимальная асимптотическая спектральная граница: большой граф не может существенно её улучшить.
        \item Таким образом, семейство \(d\)-регулярных графов Рамануджана асимптотически достигает оптимальной спектральной границы.
    \end{itemize}
    \note{
        \begin{itemize}
            \item ---
        \end{itemize}
    }
\end{frame}

% ================================================================
\begin{frame}{Любоцкий--Филлипс--Сарнак: явные графы Рамануджана}
    \begin{itemize}
        \item Берём различные простые $p\equiv q\equiv1\pmod4$, $q>2\sqrt p$;
        \item Среди целых кватернионов $a+bi+cj+dk$ нормы $a^2+b^2+c^2+d^2=p$ выбираем $a>0$ нечётное, $b,c,d$ чётные (ровно $p+1$ решений);
        \item $\exists\,t\in\mathbb{F}_q$ с $t^2=-1$ (т.к. $q\equiv1\pmod4$); каждому решению сопоставляем
        \[M=\begin{pmatrix}a+tb&c+td\\-c+td&a-tb\end{pmatrix},\quad \det M=p\pmod q;\]
        \item Берём граф Кэли матричной группы над $\mathbb F_q$:
        \begin{itemize}
            \item $\operatorname{Cay}(PSL_2(\mathbb F_q),S)$, если $p$ --- квадрат $\bmod\, q$; $\operatorname{Cay}(PGL_2(\mathbb F_q),S)$, иначе;
            \item $PGL_2(\mathbb F_q)$ --- группа обратимых матриц с точностью до ненулевого скаляра,\\ $PSL_2(\mathbb F_q)$ --- её подгруппа с определителем-квадратом;
            \item $S=\{[M]\}$ --- множество классов в $PGL_2(\mathbb F_q)$ полученных матриц;
        \end{itemize}
        \item $p+1$ выбранных матриц задают рёбра в каждой вершине, степень $d=p+1$;
        \item Получаем семейство $d$-регулярных графов, $d$ фиксировано при росте $q$;
    \end{itemize}
    \note{
        \begin{itemize}
            \item В случае, когда \(p\) не является квадратом по модулю \(q\), граф на \(PGL_2(\mathbb F_q)\) двудольный.
        \end{itemize}
    }
\end{frame}

% ================================================================
\begin{frame}{Любоцкий--Филлипс--Сарнак: явные графы Рамануджана (продолжение)}
    \begin{itemize}
        \item Оказывается, что матрица (оператор) смежности этого графа связана с функциональным оператором на пр-ве модулярных форм веса $2$;
        \item Оператор Гекке $T_p$ --- линейный оператор на МФ ($p$ --- простое);
        \item Cобственные функции и значения: $T_p f=a_p f$, $f(z)=\sum_{n=1}^{\infty}a_n e^{2\pi inz}$;
        \item Теорема (Рамануджан--Петерссон, Делинь): для нормированной параболической ($a_0=0$) собственной формы Гекке веса $k$ и уровня $N$:
        \[|a_p|\le 2p^{(k-1)/2},\ p\nmid N;\quad \text{в частности, для } k=2: |a_p|\le 2\sqrt p;\]
        % \item Теорема (г. Рамануджана--Петерссона; Делинь, 1974): если $f$ --- собственная форма оператов:
        % \[|a_p|\le 2p^{(k-1)/2};\]
        \item Откуда для нетривиальных собственных значений построенного графа:
        \[|\lambda|\le 2\sqrt p=2\sqrt{d-1},\quad d=p+1;\]
        \item Т.о., полученное семейство есть семейство графов Рамануджана, итого:
        \begin{actionenv}
            \begin{center}
            \begin{tikzpicture}[
                every node/.style={draw, rounded corners=2pt, align=center, font=\small, minimum height=9mm, inner sep=4pt}
            ]
                \node (quat) {Кватернионы};
                \node[anchor=west] (gc) at ([xshift=8mm]quat.east) {Графы Кэли};
                \node[anchor=west] (mf) at ([xshift=8mm]gc.east) {Модулярные формы};
                \node[anchor=west] (exp) at ([xshift=8mm]mf.east) {Оценка спектра};

                \draw[->] (quat) -- (gc);
                \draw[->] (gc) -- (mf);
                \draw[->] (mf) -- (exp);
            \end{tikzpicture}
            \end{center}
        \end{actionenv}
    \end{itemize}
    \note{
        \begin{itemize}
            \item ---
        \end{itemize}
    }
\end{frame}

% ================================================================
\begin{frame}{Дерево Брюа--Титса}
    \begin{itemize}
        \item С матричной группой $PGL_2(\mathbb{Q}_p)$ можно ассоциировать бесконечное $(p+1)$-регулярное дерево (дерево Брюа--Титса):
        \begin{itemize}
            \item Вершины: классы $\mathbb{Z}_p$-решёток $L\subset\mathbb{Q}_p^2$ с точностью до умножения на число из $\mathbb{Q}_p^*$;
            \item Рёбра $[L]\sim[L']$: можно выбрать представителей так, что $pL\subsetneq L'\subsetneq L$ и $[L:L']=p$;
        \end{itemize}
        \item Группа $PGL_2(\mathbb{Q}_p)$ действует на дереве; фактор по подходящей арифметической группе даёт конечный граф;
        \item Пример: $PGL_2(\mathbb Q_2)$
        \begin{center}
            \begin{tikzpicture}[
                scale=0.85,
                every node/.style={font=\footnotesize, inner sep=2pt}
            ]
                % Дерево Брюа--Титса
                \coordinate (a) at (0,1.25);
                \coordinate (b) at (-1.75,0.35);
                \coordinate (c) at (0,0.35);
                \coordinate (d) at (1.75,0.35);

                \foreach \i/\x in {1/-2.2,2/-1.3,3/-0.45,4/0.45,5/1.3,6/2.2}
                    \coordinate (t\i) at (\x,-0.55);

                % Рёбра дерева
                \draw[thick]
                    (a)--(b) (a)--(c) (a)--(d)
                    (b)--(t1) (b)--(t2)
                    (c)--(t3) (c)--(t4)
                    (d)--(t5) (d)--(t6);

                % Бесконечное продолжение
                \foreach \i in {1,...,6}{
                    \draw[dashed] (t\i)--++(-0.22,-0.48);
                    \draw[dashed] (t\i)--++(0.22,-0.48);
                }

                % Вершины
                \foreach \v in {a,b,c,d,t1,t2,t3,t4,t5,t6}
                    \fill (\v) circle[radius=2pt];

                % Метки орбит
                \node[above] at (a) {$a$};
                \node[left]  at (b) {$b$};
                \node[right] at (c) {$c$};
                \node[right] at (d) {$d$};

                \node[above left]  at (t1) {$c$};
                \node[above right] at (t2) {$d$};
                \node[above left]  at (t3) {$b$};
                \node[above right] at (t4) {$d$};
                \node[above left]  at (t5) {$b$};
                \node[above right] at (t6) {$c$};

                \node at (0,-1.3)
                    {\scriptsize Бесконечное дерево $T_2$};

                % Факторизация
                \draw[->,thick] (2.8,0)--(4.5,0)
                    node[midway,above] {\scriptsize фактор};

                % Конечный граф
                \begin{scope}[xshift=5.9cm]
                    \coordinate (A) at (0,1.1);
                    \coordinate (B) at (-0.95,-0.55);
                    \coordinate (C) at (0.95,-0.55);
                    \coordinate (D) at (0,0);

                    \draw[thick]
                        (A)--(B)--(C)--(A)
                        (A)--(D) (B)--(D) (C)--(D);

                    \foreach \v in {A,B,C,D}
                        \fill (\v) circle[radius=2pt];

                    \node[above] at (A) {$a$};
                    \node[below left] at (B) {$b$};
                    \node[below right] at (C) {$c$};
                    \node[above right] at (D) {$d$};

                    \node at (0,-1.3)
                        {\scriptsize Конечный граф $K_4$};
                \end{scope}
            \end{tikzpicture}
        \end{center}
        \item Некоторые такие арифметические факторы обладают спектральными оценками Рамануджана и образуют семейства экспандеров;
    \end{itemize}
    \note{
        \begin{itemize}
            \item ---
        \end{itemize}
    }
\end{frame}

% ================================================================
\begin{frame}{Суперсингулярные графы изогений}
    \begin{itemize}
        \item Пусть $p\ne\ell$ --- простые числа, $p\equiv1\,(12)$, $E$ --- эллиптическая кривая над $\overline{\mathbb F}_p$;
        \item Теорема: для конечной подгруппы $C\subset E(\overline{\mathbb F}_p)$ фактор $E/C$ является ЭК;
        \item При $|C|=\ell$ естественное отображение $E\to E/C$ называется $\ell$-изогенией;
        \item Суперсингулярные ЭК: $\operatorname{End}(E)$ --- макс. порядок в кватернионной алгебре;
        \item Построим граф $\mathcal{G}_p=\mathcal{G}_p(\ell)$:
        \begin{itemize}
            \item вершины --- классы изоморфизма суперсинг. ЭК над $\overline{\mathbb F}_p$ ($(p-1)/12$ вершин);
            \item рёбра $\sim$ $\ell$-изогении (возможны петли и кратные рёбра);
        \end{itemize}
        \item Теорема (Pizer, 1990): $\mathcal{G}_p(\ell)$ --- $(\ell+1)$-регулярный граф Рамануджана;
        \item Получаем семейство экспандеров $\mathcal{G}_p$, $p\to\infty$;
        \item Конструкция использует арифметику порядков и идеалов в кватернионной алгебре и оценки спектра через коэффициенты МФ; таким образом:
        \begin{actionenv}
            \begin{center}
            \begin{tikzpicture}[
                every node/.style={draw, rounded corners=2pt, align=center, font=\small, minimum height=9mm, inner sep=4pt}
            ]
                \node (ec) {Эллиптические\\кривые};
                \node[anchor=west] (quat) at ([xshift=8mm]ec.east) {Кватернионные\\алгебры};
                \node[anchor=west] (mf) at ([xshift=8mm]quat.east) {Модулярные\\формы};
                \node[anchor=west] (exp) at ([xshift=8mm]mf.east) {Графы\\экспандеры};

                \draw[<->] (ec) -- (quat);
                \draw[<->] (quat) -- (mf);
                \draw[<->] (mf) -- (exp);
            \end{tikzpicture}
            \end{center}
        \end{actionenv}
        \item Приложения: хэш-функции и постквантовая криптография.
    \end{itemize}
    \note{
        \begin{itemize}
            
            \item Условие $p\equiv1\pmod{12}$ исключает суперсингулярные ЭК с $j=0$ и $j=1728$, поскольку они суперсингулярны при $p\equiv2\pmod3$ и $p\equiv3\pmod4$ соответственно. Поэтому для всех рассматриваемых кривых $\operatorname{Aut}(E)=\{\pm1\}$ и матрица смежности графа изогений симметрична (двойственные изогении задают противоположные направления рёбер).
            \item A.~K.~Pizer, Ramanujan Graphs and Hecke Operators, Bulletin of the AMS 23 (1990), 127--137. \url{https://doi.org/10.1090/S0273-0979-1990-15918-X}
            \item A.~Pizer, Ramanujan graphs, Computational Perspectives on Number Theory, AMS, 1998, 159--178.
            \item D.~Charles, K.~Lauter, E.~Goren, Cryptographic Hash Functions from Expander Graphs, Journal of Cryptology 22 (2009), 93--113. \url{https://doi.org/10.1007/s00145-007-9002-x}
        \end{itemize}
    }
\end{frame}

% ================================================================
\begin{frame}{Хэш-функции на графах Кэли}
    \begin{itemize}
        \item Хэш-функция --- вычислимая функция, отображающая битовые строки произвольной длины в битовые строки фиксированной длины:
        \[H\colon\{0,1\}^*\to\{0,1\}^N;\]
        \item Идея конструкции: берём конечную группу $G$ и два её образующих $A_0,A_1$;
        \item Для битовой строки $w=b_1\dots b_k$ произвольной длины определяем:
        \[H(w)=A_{b_1}\cdots A_{b_k}\in G;\]
        \item Свойство: $H(uv)=H(u)H(v)$ --- можно параллелить или дописывать хэш;
        \item Интерпретация: движение по графу Кэли согласно битам сообщения;\\
        хэш --- конечная вершина маршрута;
        \item Пример: $G=SL_2(\mathbb F_p)$, где $p$ --- большое простое число,
        \[
            A_0=\left(\begin{smallmatrix}1&2\\0&1\end{smallmatrix}\right),\quad
            A_1=\left(\begin{smallmatrix}1&0\\2&1\end{smallmatrix}\right),\quad
            H(01010110)=\left(\begin{smallmatrix}309&688\\128&285\end{smallmatrix}\right)\pmod p;
        \]
        \item Теорема (Bourgain--Gamburd, 2008): графы Кэли $G=SL_2(\mathbb F_p)$ с образующими $A_0^{\pm1},A_1^{\pm1}$ образуют семейство экспандеров при $p\to\infty$;
        
    \end{itemize}
    \note{
        \begin{itemize}
            \item Хорошее расширение и большой обхват графа сами по себе не гарантируют криптографической стойкости хэш-функции.
            \item Важную роль в таких конструкциях играют свойства расширения графа и его обхват (girth) --- длина кратчайшего простого цикла.
            \item V.~Shpilrain, Girth of the Cayley graph and Cayley hash functions, 2025. \url{https://arxiv.org/abs/2502.13197}
            \item J.~Bourgain, A.~Gamburd, Uniform expansion bounds for Cayley graphs of $SL_2(\mathbb F_p)$, Annals of Mathematics 167 (2008), 625--642. \url{https://doi.org/10.4007/annals.2008.167.625}
        \end{itemize}
    }
\end{frame}

% ================================================================
\begin{frame}{Другие приложения экспандеров}
    \begin{itemize}
        \item Коды, исправляющие ошибки: экспандеры позволяют строить помехоустойчивые коды с декодированием за линейное время;
        \item Сети: построение разреженных коммуникационных сетей с хорошей связностью и без узких мест;
        \item Случайные блуждания: на экспандерах распределение положения частицы быстро приближается к равномерному независимо от начальной вершины (время перемешивания $O(\log |V|)$);
        \item Псевдослучайность и алгоритмы: заменяя независимые случайные выборки блужданием по экспандеру, можно уменьшать вероятность ошибки рандомизированных алгоритмов, расходуя меньше случайных битов;
        \item Графовые нейронные сети (GNN): экспандеры используются как дополнительная структура связей для эффективной передачи информации между удалёнными вершинами.
    \end{itemize}
    \note{
        \begin{itemize}
            \item M.~Sipser, D.~A.~Spielman, Expander Codes, IEEE Transactions on Information Theory 42 (1996), 1710--1722.
            \item S.~Hoory, N.~Linial, A.~Wigderson, Expander Graphs and Their Applications, Bulletin of the AMS 43 (2006), 439--561.
            \item A.~Deac, M.~Lackenby, P.~Veličković, Expander Graph Propagation, Learning on Graphs Conference, 2022. \url{https://arxiv.org/abs/2210.02997}
        \end{itemize}
    }
\end{frame}

% ============================================================
\section{Дзета-функция Ихары}
% ============================================================

% ================================================================
\begin{frame}{Примитивные циклы}
    \begin{itemize}
        \item Пусть $G$ --- конечный неориентированный граф;
        \item Каждому ребру сопоставим два ориентированных: $e$ и $\bar e$;
        \item Маршрут $e_1 \dots e_r$ без возвратов: $e_{i+1}\ne\bar e_i$;
        \item Далее под циклом (Ихары) будем понимать замкнутый маршрут без хвоста и без возвратов (включая $e_1\ne\bar e_r$), т.е. допускаются повторения;
        \item Цикл $C^k$ --- $k$ повторных обходов цикла $C$;
        \item Цикл $C$ --- примитивный $\leftrightharpoons$ $C\neq D^k$ для более короткого цикла $D$ и $k\ge2$;
        \item Примеры:
        \vspace{-2ex}
        \begin{center}
            \begin{actionenv}
            \begin{minipage}[t]{0.31\linewidth}
                \centering
                \begin{tikzpicture}[scale=0.95, >={Stealth[length=2mm]}]
                    \path[use as bounding box] (-1.6,-1.35) rectangle (1.6,1.35);

                    % Немедленный возврат
                    \coordinate (a1) at (-0.85,0.65);
                    \coordinate (b1) at ( 0.85,0.65);

                    \draw[thick] (a1)--(b1);
                    \draw[red,thick,->] (-0.68,0.77)--(0.68,0.77);
                    \draw[red,thick,->] (0.68,0.53)--(-0.68,0.53);

                    \foreach \v in {a1,b1}
                        \fill (\v) circle[radius=2pt];

                    \node[font=\scriptsize,red] at (0,0.3) {$e\bar e$};

                    % Замкнутый маршрут с хвостом
                    \coordinate (a) at (-1.2,-0.48);
                    \coordinate (b) at (-0.1,-0.48);
                    \coordinate (c) at ( 0.9, 0.02);
                    \coordinate (d) at ( 0.9,-0.98);

                    \draw[thick] (a)--(b)--(c)--(d)--(b);

                    % Обход треугольника
                    \foreach \u/\v in {b/c,c/d,d/b}
                        \draw[thick,->] ($(\u)!0.18!(\v)$)--($(\u)!0.78!(\v)$);

                    % Хвост
                    \draw[red,thick,->] (-1.05,-0.36)--(-0.25,-0.36);
                    \draw[red,thick,->] (-0.25,-0.60)--(-1.05,-0.60);

                    \foreach \v in {a,b,c,d}
                        \fill (\v) circle[radius=2pt];

                    \node[font=\scriptsize,red] at (-0.68,-0.94)
                        {$e_1=\bar e_5$};
                \end{tikzpicture}\\[-1ex]
                {\scriptsize Возврат и хвост запрещены}
            \end{minipage}
            \end{actionenv}
            \begin{actionenv}
            \begin{minipage}[t]{0.31\linewidth}
                \centering
                \begin{tikzpicture}[scale=0.95, >={Stealth[length=2mm]}]
                    \path[use as bounding box] (-1.8,-1.35) rectangle (1.8,1.35);

                    % Примитивный цикл C
                    \begin{scope}[xshift=-0.85cm]
                        \coordinate (a) at (-0.45, 0.45);
                        \coordinate (b) at ( 0.45, 0.45);
                        \coordinate (c) at ( 0.45,-0.45);
                        \coordinate (d) at (-0.45,-0.45);

                        \draw[thick] (a)--(b)--(c)--(d)--(a);
                        \foreach \v in {a,b,c,d}
                            \fill (\v) circle[radius=2pt];

                        \draw[->,thick] ($(a)!0.2!(b)$) -- ($(a)!0.8!(b)$);
                        \draw[->,thick] ($(b)!0.2!(c)$) -- ($(b)!0.8!(c)$);
                        \draw[->,thick] ($(c)!0.2!(d)$) -- ($(c)!0.8!(d)$);
                        \draw[->,thick] ($(d)!0.2!(a)$) -- ($(d)!0.8!(a)$);

                        \node[font=\scriptsize] at (0,-0.9) {$C$};
                    \end{scope}

                    % Непримитивный цикл C^2
                    \begin{scope}[xshift=0.85cm]
                        \coordinate (a) at (-0.45, 0.45);
                        \coordinate (b) at ( 0.45, 0.45);
                        \coordinate (c) at ( 0.45,-0.45);
                        \coordinate (d) at (-0.45,-0.45);

                        \draw[thick] (a)--(b)--(c)--(d)--(a);
                        \foreach \v in {a,b,c,d}
                            \fill (\v) circle[radius=2pt];

                        \draw[->,thick] ($(a)!0.15!(b)$) -- ($(a)!0.75!(b)$);
                        \draw[->,thick] ($(b)!0.15!(c)$) -- ($(b)!0.75!(c)$);
                        \draw[->,thick] ($(c)!0.15!(d)$) -- ($(c)!0.75!(d)$);
                        \draw[->,thick] ($(d)!0.15!(a)$) -- ($(d)!0.75!(a)$);

                        \draw[->,thick] (-0.28,0.28) -- (0.28,0.28);
                        \draw[->,thick] (0.28,0.28) -- (0.28,-0.28);
                        \draw[->,thick] (0.28,-0.28) -- (-0.28,-0.28);
                        \draw[->,thick] (-0.28,-0.28) -- (-0.28,0.28);

                        \node[font=\scriptsize] at (0,0) {$2\times$};
                        \node[font=\scriptsize] at (0,-0.9) {$C^2$};
                    \end{scope}
                \end{tikzpicture}\\[-1ex]
                {\scriptsize $C$ примитивен, $C^2$ --- нет}
            \end{minipage}
            \end{actionenv}
            \begin{actionenv}
            \begin{minipage}[t]{0.31\linewidth}
                \centering
                \begin{tikzpicture}[scale=0.95, >={Stealth[length=2mm]}]
                    \path[use as bounding box] (-1.65,-1.35) rectangle (1.65,1.35);

                    % "Бабочка"
                    \coordinate (o)  at (0,0);
                    \coordinate (a)  at (-1.0, 0.55);
                    \coordinate (b)  at (-1.0,-0.55);
                    \coordinate (c)  at ( 1.0, 0.55);
                    \coordinate (d)  at ( 1.0,-0.55);

                    \draw[thick] (o)--(a)--(b)--(o)--(c)--(d)--(o);

                    \foreach \v in {o,a,b,c,d}
                        \fill (\v) circle[radius=2pt];

                    % Ориентированный примитивный цикл
                    \draw[->,thick] ($(o)!0.2!(a)$) -- ($(o)!0.8!(a)$);
                    \draw[->,thick] ($(a)!0.2!(b)$) -- ($(a)!0.8!(b)$);
                    \draw[->,thick] ($(b)!0.2!(o)$) -- ($(b)!0.8!(o)$);
                    \draw[->,thick] ($(o)!0.2!(c)$) -- ($(o)!0.8!(c)$);
                    \draw[->,thick] ($(c)!0.2!(d)$) -- ($(c)!0.8!(d)$);
                    \draw[->,thick] ($(d)!0.2!(o)$) -- ($(d)!0.8!(o)$);

                \end{tikzpicture}\\[-1ex]
                {\scriptsize Примитивный, не простой}
            \end{minipage}
            \end{actionenv}
        \end{center}
    \end{itemize}
    \note{
        \begin{itemize}
            \item ---
        \end{itemize}
    }
\end{frame}

% ================================================================
\begin{frame}{Дзета-функция Ихары}
    \begin{itemize}
        \item Циклы отождествляются при циклическом сдвиге начальной точки;
        \item Обозначим через $[C]$ класс примитивного цикла;
        \item $\ell([C])=\ell(C)$ --- длина примитивного цикла (число рёбер при одном обходе);
        \item Дзета-функция --- эйлерово произведение по всем таким классам:
        \[Z_G(u)=\prod_{[C]}\left(1-u^{\ell(C)}\right)^{-1};\]
        \item Другие дзета-функции:
        \begin{itemize}
            \item $\zeta(s)=\prod_p (1-p^{-s})^{-1}$\quad --- Римана;
            \item $\zeta_{\mathbb{Z}[i]}(s)=\prod_{\mathfrak p}(1-N(\mathfrak p)^{-s})^{-1}$\quad --- кольца гауссовых целых;
            \item $\zeta_{\mathbb{F}_q[x]}(s)=\prod_f (1-q^{-s\deg f})^{-1}$\quad --- кольца многочленов над $\mathbb{F}_q$;
        \end{itemize}
        \item Аналогия: простые идеалы $\leftrightarrow$ классы примитивных циклов.
    \end{itemize}
    \note{
        \begin{itemize}
            \item ---
        \end{itemize}
    }
\end{frame}

% ================================================================
\begin{frame}{Дзета-функция и подсчёт маршрутов}
    \begin{itemize}
        \item Пример: дерево --- нет циклов: $Z_T(u)=1$;
        \item Пример: граф-цикла --- два примитивных класса $\ell(n)=n$ (две ориентации):
        \[Z_{C_n}(u)=\frac{1}{(1-u^n)^2};\]
        \item В общем случае произведение $\prod_{[C]}$ может быть бесконечным;
        \item Вспомним также локальную дзета-функцию (эллиптической) кривой:
        \[\log \zeta_{E,p}(u) = \sum_{m=1}^\infty \frac{N_{p,m}(E)}{m} u^m;\]
        \item Пусть $N_m$ --- число циклов длины $m$ (не обязательно примитивных),\\ с отмеченным первым ребром;
        \item Тогда:
        \[\log Z_G(u)=\sum_{m\ge1}\frac{N_m}{m}u^m;\]
    \end{itemize}
    \note{
        \begin{itemize}
            \item ---
        \end{itemize}
    }
\end{frame}

% ================================================================
\begin{frame}{Рациональность дзета-функции}
    \begin{itemize}
        \item Заменим каждое ребро графа двумя ориентированными: $e$ и $\bar e$;
        \item Рассмотрим матрицу $B$ размера $2|E|\times2|E|$:
        \[b_{ef}=\begin{cases}1,&t(e)=o(f)\text{ и }f\ne\bar e,\\0,&\text{иначе};\end{cases}\]
        где $o(e)$, $t(e)$ --- начальная и конечная вершины ориентированного ребра;
        \item Лемма: $\operatorname{tr}(B^m)=N_m$;
        \item Теорема: $Z_G(u)^{-1}=\det(I-uB)$;
        \item Таким образом, $Z_G(u)$ --- рациональная функция;
        \item Пусть далее $G$ --- связный граф, $n=|V|$, $m=|E|$;
        \item $A$ --- матрица смежности, $D=\operatorname{diag}(\deg(v))$;
        \item Теорема (Ихара--Басс): $Z_G(u)^{-1}=(1-u^2)^{m-n}\det\left(I-uA+u^2(D-I)\right)$;
    \end{itemize}
    \note{
        \begin{itemize}
            \item ---
        \end{itemize}
    }
\end{frame}

% ================================================================
\begin{frame}{Дзета-функция регулярных графов}
    \begin{itemize}
        \item Если далее граф $d$-регулярный ($D=dI$), $d\ge3$, $\lambda_1,\dots,\lambda_n$ --- спектр, то:
        \[Z_G(u)^{-1}=(1-u^2)^{m-n}\prod_{i=1}^n\bigl(1-\lambda_i u+(d-1)u^2\bigr);\]
        \item Тривиальные нули $Z_G(u)^{-1}$ --- нули от $(1-u^2)^{m-n}$ и собственных значений
        $\lambda=\pm d$ (если они есть): $u=\pm1$, $u=\pm\frac{1}{d-1}$;
        \item Аналог г. Римана: $\forall$ нетривиального нуля $|u|=1/\sqrt{d-1}$;
        \item Нетривиальному собственному значению $\lambda$ соответствует уравнение
        \[1-\lambda u+(d-1)u^2=0;\]

        \item Его корни лежат на окружности $|u|=1/\sqrt{d-1}$ $\Leftrightarrow$ $|\lambda|\le2\sqrt{d-1}$,
        \item Это верно для всех нетривиальных $\lambda$ $\Leftrightarrow$ $G$ --- граф Рамануджана;
        
        \item Теорема: для конечного связного $d$-регулярного графа при $d\ge3$:\\
        $G$ --- граф Рамануджана $\Leftrightarrow$ $Z_G(u)^{-1}$ удовлетворяет аналогу г. Римана.
    \end{itemize}
    \note{
        \begin{itemize}
            \item ---
        \end{itemize}
    }
\end{frame}

% ================================================================
\begin{frame}{Приложения дзета-функций графов}
    \begin{itemize}
        \item Примитивные циклы: существует аналог теоремы о распределении простых чисел, описывающий асимптотику числа таких циклов заданной длины;
        \item Остовные деревья: при $|E|>|V|$ их число выражается через производную определителя из формулы Ихары--Басса при $u=1$;
        \item Спектр: для $d$-регулярных графов с фиксированными $d$ и $|V|$ дзета-функция определяет спектр матрицы смежности;
        \item Инварианты графов: существуют неизоморфные графы с одинаковой дзета-функцией Ихары;
        \item Теория Галуа: для накрытий графов $Y\to X$ строится аналог теории Галуа расширений полей, включая законы расщепления примитивных циклов;
        \item Факторизация Артина: дзета-функция $Z_Y$ регулярного накрытия раскладывается в произведение графовых $L$-функций Артина.
    \end{itemize}
    \note{
        \begin{itemize}
            \item Накрытия графов --- топологические накрытия графов как одномерных CW-комплексов. Для связного регулярного накрытия группа дековых автоморфизмов играет роль группы Галуа.
            \item S.~Northshield, A Note on the Zeta Function of a Graph, Journal of Combinatorial Theory, Series B 74 (1998), 408--410.
            \item H.~M.~Stark, A.~A.~Terras, Zeta Functions of Finite Graphs and Coverings, Part II, Advances in Mathematics 154 (2000), 132--195.
            \item A.~Terras, Zeta Functions of Graphs: A Stroll through the Garden, Cambridge Univ. Press, 2010, гл.~10, 13--22.
        \end{itemize}
    }
\end{frame}

% ============================================================
\section{Кубик Рубика, графы, RL}
% ============================================================

% ================================================================
\begin{frame}{Кубик Рубика как группа}
    \begin{itemize}
        \item Допустимые преобразования кубика Рубика $3\times3\times3$ образуют группу $G$;
        \item Группа порождается шестью поворотами граней на $\pi/2$: $G=\langle U,D,L,R,F,B\rangle$;
        \item Пример:
        \vspace{-1ex}
        \begin{center}
        % \begin{tikzpicture}[line join=round]
        %     \def\cubestep{2.9}
        %     \colorlet{rubG}{green!65!black}
        %     \colorlet{rubR}{red!85!black}
        %     \colorlet{rubB}{blue!75!black}
        %     \colorlet{rubY}{yellow!90!orange}
        %     \colorlet{rubW}{white}
        %     \tikzset{sticker/.style={draw=black,line width=0.7pt}}

        %     \foreach \st in {0,1,2,3} {
        %         \ifcase\st\relax
        %             % e
        %             \def\frontcolors{rubG,rubG,rubG,rubG,rubG,rubG,rubG,rubG,rubG}
        %             \def\rightcolors{rubR,rubR,rubR,rubR,rubR,rubR,rubR,rubR,rubR}
        %             \def\topcolors{rubW,rubW,rubW,rubW,rubW,rubW,rubW,rubW,rubW}
        %             \def\word{e}
        %         \or
        %             % U
        %             \def\frontcolors{rubG,rubG,rubG,rubG,rubG,rubG,rubR,rubR,rubR}
        %             \def\rightcolors{rubR,rubR,rubR,rubR,rubR,rubR,rubB,rubB,rubB}
        %             \def\topcolors{rubW,rubW,rubW,rubW,rubW,rubW,rubW,rubW,rubW}
        %             \def\word{U}
        %         \or
        %             % UR
        %             \def\frontcolors{rubG,rubG,rubY,rubG,rubG,rubY,rubR,rubR,rubY}
        %             \def\rightcolors{rubR,rubR,rubB,rubR,rubR,rubB,rubR,rubR,rubB}
        %             \def\topcolors{rubW,rubW,rubG,rubW,rubW,rubG,rubW,rubW,rubR}
        %             \def\word{UR}
        %         \or
        %             % URU^{-1}
        %             \def\frontcolors{rubG,rubG,rubY,rubG,rubG,rubY,rubG,rubG,rubG}
        %             \def\rightcolors{rubR,rubR,rubB,rubR,rubR,rubB,rubR,rubR,rubY}
        %             \def\topcolors{rubW,rubW,rubW,rubW,rubW,rubW,rubR,rubG,rubG}
        %             \def\word{URU^{-1}}
        %         \fi

        %         \pgfmathsetmacro{\dx}{\cubestep*\st}
        %         \begin{scope}[xshift=\dx cm]
        %             \begin{scope}[x={(0.28cm,-0.16cm)},y={(0.28cm,0.16cm)},z={(0cm,0.32cm)}]
        %                 % Передняя грань
        %                 \foreach \clr [count=\k from 0] in \frontcolors {
        %                     \pgfmathtruncatemacro{\i}{mod(\k,3)}
        %                     \pgfmathtruncatemacro{\j}{int(\k/3)}
        %                     \filldraw[sticker,fill=\clr]
        %                         (\i,0,\j) -- ({\i+1},0,\j) --
        %                         ({\i+1},0,{\j+1}) -- (\i,0,{\j+1}) -- cycle;
        %                 }
        %                 % Правая грань
        %                 \foreach \clr [count=\k from 0] in \rightcolors {
        %                     \pgfmathtruncatemacro{\i}{mod(\k,3)}
        %                     \pgfmathtruncatemacro{\j}{int(\k/3)}
        %                     \filldraw[sticker,fill=\clr]
        %                         (3,\i,\j) -- (3,{\i+1},\j) --
        %                         (3,{\i+1},{\j+1}) -- (3,\i,{\j+1}) -- cycle;
        %                 }
        %                 % Верхняя грань
        %                 \foreach \clr [count=\k from 0] in \topcolors {
        %                     \pgfmathtruncatemacro{\i}{mod(\k,3)}
        %                     \pgfmathtruncatemacro{\j}{int(\k/3)}
        %                     \filldraw[sticker,fill=\clr]
        %                         (\i,\j,3) -- ({\i+1},\j,3) --
        %                         ({\i+1},{\j+1},3) -- (\i,{\j+1},3) -- cycle;
        %                 }
        %             \end{scope}

        %             % Подпись состояния
        %             \node[anchor=base] at (0.84,-0.85) {$\word$};
        %         \end{scope}
        %     }

        %     % Стрелки и подписи действий
        %     \foreach \st/\move in {0/U,1/R,2/{U^{-1}}} {
        %         \pgfmathsetmacro{\xa}{\cubestep*\st+1.80}
        %         \pgfmathsetmacro{\xb}{\cubestep*(\st+1)-0.12}
        %         \draw[->,thick] (\xa,0.52) -- (\xb,0.52);
        %         \node[anchor=base] at ({(\xa+\xb)/2},0.67) {$\move$};
        %     }
        % \end{tikzpicture}

\begin{tikzpicture}[line join=round]
    \def\cubestep{2.9}
    \colorlet{rubG}{green!65!black}
    \colorlet{rubR}{red!85!black}
    \colorlet{rubB}{blue!75!black}
    \colorlet{rubY}{yellow!90!orange}
    \colorlet{rubW}{white}
    \tikzset{sticker/.style={draw=black,line width=0.7pt}}

    % Стили анимации
    \ifanimated
        \tikzset{
            cube step/.style={visible on=<+->},
            cube label/.style={alt=<.>{text=ThemeAlert}{}},
            cube arrow/.style={alt=<.>{draw=ThemeAlert}{}}
        }
    \else
        \tikzset{
            cube step/.style={},
            cube label/.style={},
            cube arrow/.style={}
        }
    \fi

    \foreach \st in {0,1,2,3} {
        \ifcase\st\relax
            % e
            \def\frontcolors{rubG,rubG,rubG,rubG,rubG,rubG,rubG,rubG,rubG}
            \def\rightcolors{rubR,rubR,rubR,rubR,rubR,rubR,rubR,rubR,rubR}
            \def\topcolors{rubW,rubW,rubW,rubW,rubW,rubW,rubW,rubW,rubW}
            \def\word{e}
        \or
            % U
            \def\frontcolors{rubG,rubG,rubG,rubG,rubG,rubG,rubR,rubR,rubR}
            \def\rightcolors{rubR,rubR,rubR,rubR,rubR,rubR,rubB,rubB,rubB}
            \def\topcolors{rubW,rubW,rubW,rubW,rubW,rubW,rubW,rubW,rubW}
            \def\word{U}
            \def\move{U}
        \or
            % UR
            \def\frontcolors{rubG,rubG,rubY,rubG,rubG,rubY,rubR,rubR,rubY}
            \def\rightcolors{rubR,rubR,rubB,rubR,rubR,rubB,rubR,rubR,rubB}
            \def\topcolors{rubW,rubW,rubG,rubW,rubW,rubG,rubW,rubW,rubR}
            \def\word{UR}
            \def\move{R}
        \or
            % URU^{-1}
            \def\frontcolors{rubG,rubG,rubY,rubG,rubG,rubY,rubG,rubG,rubG}
            \def\rightcolors{rubR,rubR,rubB,rubR,rubR,rubB,rubR,rubR,rubY}
            \def\topcolors{rubW,rubW,rubW,rubW,rubW,rubW,rubR,rubG,rubG}
            \def\word{URU^{-1}}
            \def\move{U^{-1}}
        \fi

        % Первый кубик появляется вместе с буллитом
        \ifanimated
            \ifnum\st=0\relax
                \tikzset{cube step/.style={visible on=<.->}}
            \fi
        \fi

        \pgfmathsetmacro{\dx}{\cubestep*\st}
        \begin{scope}[xshift=\dx cm,cube step]
            \begin{scope}[
                x={(0.28cm,-0.16cm)},
                y={(0.28cm,0.16cm)},
                z={(0cm,0.32cm)}
            ]
                % Передняя грань
                \foreach \clr [count=\k from 0] in \frontcolors {
                    \pgfmathtruncatemacro{\i}{mod(\k,3)}
                    \pgfmathtruncatemacro{\j}{int(\k/3)}
                    \filldraw[sticker,fill=\clr]
                        (\i,0,\j) -- ({\i+1},0,\j) --
                        ({\i+1},0,{\j+1}) -- (\i,0,{\j+1}) -- cycle;
                }

                % Правая грань
                \foreach \clr [count=\k from 0] in \rightcolors {
                    \pgfmathtruncatemacro{\i}{mod(\k,3)}
                    \pgfmathtruncatemacro{\j}{int(\k/3)}
                    \filldraw[sticker,fill=\clr]
                        (3,\i,\j) -- (3,{\i+1},\j) --
                        (3,{\i+1},{\j+1}) -- (3,\i,{\j+1}) -- cycle;
                }

                % Верхняя грань
                \foreach \clr [count=\k from 0] in \topcolors {
                    \pgfmathtruncatemacro{\i}{mod(\k,3)}
                    \pgfmathtruncatemacro{\j}{int(\k/3)}
                    \filldraw[sticker,fill=\clr]
                        (\i,\j,3) -- ({\i+1},\j,3) --
                        ({\i+1},{\j+1},3) -- (\i,{\j+1},3) -- cycle;
                }
            \end{scope}

            % Подпись состояния
            \node[anchor=base,cube label]
                at (0.84,-0.85) {$\word$};

            % Стрелка от предыдущего состояния
            \ifnum\st>0\relax
                \draw[->,thick,cube arrow]
                    ({1.80-\cubestep},0.52) -- (-0.12,0.52);

                \node[anchor=base,cube label]
                    at ({(1.68-\cubestep)/2},0.67) {$\move$};
            \fi
        \end{scope}
    }
\end{tikzpicture}

        \end{center}
        \item Элементы $g\in G$ отождествляются с достижимыми состояниями кубика относительно собранного ($e\in G$);
        \item Число достижимых состояний:
        \[|G|=43\,252\,003\,274\,489\,856\,000\approx4{,}3\cdot10^{19};\]
    \end{itemize}
    \note{
        \begin{itemize}
            \item ---
        \end{itemize}
    }
\end{frame}

% ================================================================
\begin{frame}{Граф Кэли кубика Рубика}
    \begin{itemize}
        \item Пусть $S$ --- множество допустимых элементарных ходов, $S=S^{-1}$, $\langle S\rangle=G$;
        \item Например, способы считать элементраные ходы:
        \begin{itemize}
            \item FTM (face-turn metric): повороты граней на $\pm\pi/2$ и $\pi$ ($|S|=18$);
            \item QTM (quarter-turn metric): только повороты граней на $\pm\pi/2$ ($|S|=12$);
        \end{itemize}
        \item Рассмотрим соответствующий граф Кэли $\operatorname{Cay}(G,S)$;
        \item Решить кубик --- найти путь в $\operatorname{Cay}(G,S)$ от состояния $g$ до единицы $e$; длина кратчайшего решения --- расстояние в графе $d_S(g,e)$;
        \item Число Бога (God's number) --- максимальная длина кратчайшего решения, то есть диаметр графа Кэли: $\operatorname{diam}_S=\operatorname{diam} \operatorname{Cay}(G,S)$;
        \item Теорема (компьютерное доказательство):
        \begin{itemize}
            \item $\operatorname{diam}_{\mathrm{FTM}}=20$ (Rokicki, Kociemba, Davidson, Dethridge, 2010);
            \item $\operatorname{diam}_{\mathrm{QTM}}=26$ (Rokicki, Davidson, 2014).
        \end{itemize}
        \item Для конкретного решения поиск в ширину (BFS) находит кратчайший путь, но для огромного графа требует слишком много памяти и времени;
    \end{itemize}
    \note{
    \begin{itemize}
        \item Оба результата получены с помощью компьютерных доказательств; некомпьютерные доказательства точных значений диаметра неизвестны.
        \item Нижние оценки $20$ для FTM и $26$ для QTM были установлены Майклом Ридом в 1995 и 1998 гг. соответственно.
        \item Пространство состояний разбито на $2\,217\,093\,120$ смежных классов подгруппы $H=\langle U,D,F^2,B^2,L^2,R^2\rangle$. Использование симметрий и решение задачи о покрытии позволили сократить число обрабатываемых классов до $55\,882\,296$.
        \item Вычисления для FTM (2010) заняли около 35 процессорных лет на ресурсах Google, для QTM (2014) --- около 29 процессорных лет в Ohio Supercomputer Center.
        \item Исходный код доказательства для FTM опубликован на CWEB (около 5100 строк в четырёх основных модулях, включая документацию). Корректность проверялась сравнением независимых реализаций алгоритмов и независимой проверкой покрытия смежных классов.
        \item Источник: Rokicki, Kociemba, Davidson, Dethridge, The Diameter of the Rubik's Cube Group Is Twenty, SIAM Journal on Discrete Mathematics 27 (2013), 1082--1105.
        \item \url{https://www.cube20.org/}, \url{https://www.cube20.org/qtm/}, \url{https://www.cube20.org/src/}.
    \end{itemize}
    }
\end{frame}

% ================================================================
\begin{frame}{Обучение с подкреплением (Reinforcement Learning)}
    \begin{itemize}
        \item Обучение с подкреплением (RL) --- подход к ML, при котором агент учится выбирать действия, взаимодействуя со средой и получая награды;
        \item Цель обучения --- найти стратегию действий (политику), максимизирующую ожидаемую суммарную награду;
        \item В случае кубика Рубика:
        \begin{itemize}
            \item состояния --- элементы группы $g\in G$;
            \item действия --- элементарные ходы $s\in S$;
            \item переходы между состояниями: $g\mapsto gs$;
            \item награда за каждый ход равна $-1$; при достижении $e$ процесс завершается;
        \end{itemize}
        \item Максимизация суммарной награды $\sim$ минимизации числа ходов до цели;
        \item Оптимальная функция стоимости $J^*(g)=d_S(g,e)$ удовлетворяет уравнению Беллмана:
        \vspace{-1ex}
        \[J^*(g)=1+\min_{s\in S}J^*(gs),\quad g\ne e;\qquad J^*(e)=0;\]
        \item Идея: вместо вычисления всех $d_S(g,e)$ научиться приближать функцию $J^*$.
    \end{itemize}
    \note{
        \begin{itemize}
            \item ---
        \end{itemize}
    }
\end{frame}

% ================================================================
\begin{frame}{DeepCubeA: обучение на графе Кэли}
    \begin{itemize}
        \item Алгоритм DeepCubeA (Agostinelli и др., 2019): RL + эвристический поиск;
        \item Обучающие состояния генерируются в обратном направлении: от собранного кубика случайными последовательностями ходов;
        \item Нейросеть $J_\theta(g)$ оценивает оставшееся число ходов; для обучения используются её собственные оценки соседних состояний:
        \[J_{\mathrm{target}}(g)=1+\min_{s\in S}J_\theta(gs),\quad g\ne e,\qquad J_{\mathrm{target}}(e)=0;\]
        \item Информация о расстоянии до цели постепенно распространяется по графу;
        \item После обучения оценка $J_\theta$ используется в модификации алгоритма A* --- поиска на графе, при котором вершины выбираются с учётом длины уже пройденного пути и оценки оставшегося расстояния;
        \item Результат: обученная эвристика позволяет исследовать меньше вершин, быстрее находить решения и получать пути, близкие к кратчайшим, без использования специальных знаний о группе кубика.
        % \item Результат: по сравнению с классической эвристикой Корфа обученная эвристика в экспериментах позволяет генерировать меньше вершин, быстрее находить решения и получать пути, близкие к кратчайшим, без использования специальных знаний о группе кубика.
    \end{itemize}
    \note{
        \begin{itemize}
            \item На тестовой выборке из $1000$ конфигураций (QTM) решены все; $60{,}3\%$ решений оптимальны, остальные длиннее кратчайших не более чем на четыре хода.
            \item Agostinelli F., McAleer S., Shmakov A., Baldi P. Solving the Rubik's cube with deep reinforcement learning and search. Nature Machine Intelligence 1 (2019), 356--363.
            \item \url{https://doi.org/10.1038/s42256-019-0070-z}
            \item Интерактивная демонстрация алгоритма: \url{https://deepcube.igb.uci.edu/}
        \end{itemize}
    }
\end{frame}

% ================================================================
\begin{frame}{Спектральные свойства графов и RL}
    \begin{itemize}
        \item Задача в RL: исследование пр-ва состояний (exploration) с редкой наградой;
        \item Пример: робот исследует лабиринт, не зная, где находится цель;\\ случайные действия задают блуждание по графу состояний;
        \item Время покрытия $T_{\mathrm{cov}}(G)$ --- ожидаемое число шагов случайного блуждания до посещения всех вершин в худшем случае по начальной вершине;
        \item Пример: для циклов $T_{\mathrm{cov}}(C_n)=\frac{n(n-1)}{2}=\Theta(n^2)$\quad (\,$f=\Theta(g)$ $\leftrightharpoons$ $c_1 g < f < c_2 g$);
        \item Теорема: для экспандеров $T_{\mathrm{cov}}(G_n)=\Theta(n\log n)$\quad ($n$ вершин, фикс. степень);\\
        экспандеры достигают асимптотически оптимального порядка $T_{\mathrm{cov}}$;
        \item Идея: расширить множество действий агента составными действиями,\\ т.е. добавить новые рёбра в граф состояний;
        \item Jinnai и др. (2019, 2020): выбор таких действий через спектральную теорию\\
        $\to$ уменьшение верхней оценки $T_{\mathrm{cov}}(G)$\\
        $\to$ ускорение исследования среды и всего процесса обучения.
    \end{itemize}
    
    \note{
        \begin{itemize}
            \item Нижняя оценка времени покрытия произвольного связного неориентированного графа: $T_{\mathrm{cov}}(G)\ge(1+o(1))n\log n$ (Feige, 1995).
            \item Оценка для экспанедров: Broder A. Z., Karlin A. R. Bounds on the cover time. Journal of Theoretical Probability 2 (1989), 101--120. \url{https://doi.org/10.1007/BF01048273}
            \item Jinnai et al., Discovering Options for Exploration by Minimizing Cover Time, ICML 2019.
            \item Jinnai et al., Exploration in Reinforcement Learning with Deep Covering Options, ICLR 2020.
        \end{itemize}
    }
\end{frame}

% ================================================================
\begin{frame}{Итоги лекции}
    \begin{itemize}
        \item Арифметические объекты дают естественные графы: базисы решёток, $p$-адические приближения, группы, изогении ЭК;
        \item Графы экспандеры --- графы с «хорошими и удобными» свойствами связности; их свойства связаны со спектром матрицы смежности;
        \item Кватернионы, конечные группы, деревья Брюа--Титса и модулярные формы приводят к явным конструкциям графов Рамануджана;
        \item Примитивные циклы графа играют роль «простых»; формула Ихары--Басса связывает дзета-функцию графа $Z_G(u)$ и его матрицу смежности;
        \item Для связного $d$-регулярного графа ($d\ge3$): $G$ --- граф Рамануджана $\Leftrightarrow$ выполнен аналог гипотезы Римана для $Z_G(u)$;
        \item Спектральные свойства графов и экспандеры имеют многочисленные приложения.
    \end{itemize}
    \note{
        \begin{itemize}
            \item ---
        \end{itemize}
    }
\end{frame}

% ================================================================
\begin{frame}{Литература}
    \begin{itemize}[<.->]
        \item G.~Davidoff, P.~Sarnak, A.~Valette, \textit{Elementary Number Theory, Group Theory and Ramanujan Graphs}, Cambridge Univ. Press, 2003;
        \item A.~Terras, \textit{Zeta Functions of Graphs: A Stroll through the Garden}, Cambridge Univ. Press, 2010;
        \item P.~Sarnak, \textit{Some Applications of Modular Forms}, Cambridge Univ. Press, 1990.
    \end{itemize}
    \note{
        \begin{itemize}
            \item ---
        \end{itemize}
    }
\end{frame}

\end{document}
